% \chapter{Введение}

% При анализе данных классические методы построения оценок положения и разброса, такие как выборочное среднее и ковариационная матрица, часто опираются на определенные предположения о распределении данных. Однако, реальные выборки часто содержат выбросы или не соответствуют этим предположениям, что может существенно искажать результаты и приводить к неверным выводам. В таких случаях робастные методы становятся более предпочтительными, поскольку они обладают меньшей чувствительностью к выбросам и потому устойчивы.

% В рамках предыдущей работы~\cite{filatova2024robust} был предложен и исследован алгоритм для оценки многомерного параметра положения, который является робастным.  В ходе исследования были изучены его основные свойства, включая сходимость, временную сложность и робастность на различных данных, а также было проведено сравнение с методами оценки параметра положения, такими как медиана Тьюки~\cite{tukey1975mathematics}, Оджа~\cite{oja1983descriptive} и геометрическая медиана~\cite{weber1929theory}.

% Дальнейшее исследование этого алгоритма является актуальным, поскольку позволяет не только более детально изучить его свойства и модификации, повышающие точность и устойчивость, но и оценить параметр разброса. 

% Настоящая работа является продолжением этого исследования. Основное внимание уделяется модификациям алгоритма и разработке робастной оценки матрицы разброса. Также исследуется его поведение в условиях, когда алгоритм не сходится, и проводится эмпирическая оценка его эффективности для различных типов распределений и областей. 

% \vspace{1.7ex}

% \textbf{Цель выпускной работы:}

%  Углубленное исследование предложенного ранее робастного алгоритма для оценки многомерного параметра положения, с фокусом на изучение его сходимости при работе с различными типами областей и распределений и повышение его эффективности с помощью исследования модификаций, а также разработку и изучение робастной оценки матрицы разброса.

% \vspace{2ex}

% \textbf{Задачи  выпускной работы:}

% Для достижения поставленной цели были определены следующие задачи:
% \begin{itemize}
%     \item Ознакомиться с современной теорией робастного оценивания и ее предположениями.
%     \item Предложить и проанализировать модификации предложенного ранее алгоритма оценки параметра положения, провести его теоретический анализ.
%     \item Предложить и проанализировать робастный метод оценивания матрицы разброса.
%     \item Провести анализ работы алгоритма для различных типов областей (несвязных, выпуклых, невыпуклых) и дискретных распределений.
% \end{itemize}


\chapter{Введение}

Классические методы оценивания параметров многомерных распределений, такие как выборочное среднее и выборочная ковариационная матрица, широко используются при анализе данных. Однако данные методы существенно чувствительны к выбросам и отклонениям от предполагаемой модели распределения. Даже небольшое количество аномальных наблюдений может приводить к значительным искажениям оценок и, как следствие, к неверным выводам при анализе данных.

В связи с этим важную роль играют робастные методы оценивания, обладающие устойчивостью к выбросам и нарушениям предположений о распределении данных. Особый интерес представляет задача робастного оценивания многомерного параметра положения и матрицы разброса, поскольку в многомерном случае отсутствует естественное обобщение одномерной медианы, а многие известные робастные оценки оказываются вычислительно сложными.

В рамках предыдущей работы~\cite{filatova2024robust} был предложен стохастический алгоритм робастного оценивания многомерного параметра положения. В ходе исследования были изучены его основные свойства, включая поведение алгоритма на различных типах распределений, вычислительную сложность и робастность. Кроме того, проводилось сравнение предложенного подхода с известными методами многомерного робастного оценивания, такими как медиана Тьюки~\cite{tukey1975mathematics}, медиана Оджа~\cite{oja1983descriptive} и геометрическая медиана~\cite{weber1929theory}.

Дальнейшее исследование данного алгоритма представляет интерес по нес\-кольким причинам. Во-первых, требуется более детальное изучение его теоретических свойств и связи с известными робастными оценками многомерного параметра положения. Во-вторых, стохастическая природа алгоритма приводит к необходимости исследования модификаций, повышающих устойчивость и точность оценивания. Наконец, естественным продолжением данного подхода является построение робастной оценки матрицы разброса.

Настоящая работа посвящена дальнейшему исследованию предложенного стохастического алгоритма, изучению его теоретических свойств, а также разработке модификаций, направленных на повышение устойчивости и точности оценивания. Особое внимание уделяется анализу поведения алгоритма в случаях отсутствия устойчивой сходимости и построению робастных методов оценивания матрицы разброса.

\vspace{1.7ex}

\textbf{Цель выпускной работы}

Исследование свойств предложенного ранее стохастического робастного алгоритма оценивания многомерного параметра положения, разработка и анализ его модификаций, направленных на повышение устойчивости и точности, а также построение робастных методов оценивания матрицы разброса.

\vspace{2.2ex}

\textbf{Задачи выпускной работы}

Для достижения поставленной цели были определены следующие задачи:
\begin{itemize}
    \item Изучить современные подходы к робастному оцениванию многомерного параметра положения и матрицы разброса;
    
    \item Исследовать теоретические свойства предложенного стохастического алгоритма и установить его связь с известными робастными оценками;
    
    \item Разработать и проанализировать модификации алгоритма, направленные на повышение устойчивости и точности оценивания;
    
    \item Исследовать поведение алгоритма на выборках из различных типов распределений;
    
    \item Исследовать случаи отсутствия устойчивой сходимости алгоритма;
    
    \item Разработать и исследовать робастные методы оценивания матрицы разброса.
\end{itemize}